% ===========================================================================
%  Ficha XGodel — Operadores Matemáticos · Primero de Secundaria
%  Compilar:  xelatex ficha-operadores-matematicos.tex
% ===========================================================================
\documentclass{xgodel-ficha}

\marca{XGodel}
\sitio{www.XGodel.com}
\grado{Primero de Secundaria}
\curso{Raz. Matemático}
\titulo{Operadores Matemáticos}

\begin{document}
\portada

\begin{ficha}

% --------------------------- COLUMNA IZQUIERDA ---------------------------
\begin{pensamiento}
Hoy veremos el capítulo más sencillo del 1\textsuperscript{er} Bimestre,
sólo tienes que tener en cuenta las 4 operaciones fundamentales.
\end{pensamiento}

\concepto{Operación Matemática}
\lineas{3}

\concepto{Operador Matemático}
\lineas{3}

\begin{center}
\renewcommand{\arraystretch}{1.35}
\begin{tabular}{|>{\raggedright\arraybackslash}p{25mm}|>{\centering\arraybackslash}p{21mm}|}
\hline
\textbf{OPERACIÓN} & \textbf{OPERADOR} \\ \hline
Suma            & $+$      \\
Resta           & $-$      \\
Multiplicación  & $\times$ \\
División        & $\div$   \\
\multicolumn{1}{|c|}{$\vdots$} & \\ \hline
\end{tabular}
\end{center}

\begin{globo}
Nosotros podemos definir nuevas \textbf{Operaciones Matemáticas} con las ya
existentes usando \textbf{Nuevos Operadores}, como se verá ahora.
\end{globo}

% ---------------------- FORMAS DE DEFINIR UNA OPERACIÓN -------------------
\begin{apartados}

\item \rotulo{Mediante Fórmula}
  Ejemplo:
  \formula{$a \op b \;=\; 2a + 3b$}
  Luego:
  \par\smallskip
  $1 \op 2 =$ \par\smallskip
  $3 \op 5 =$

  \par\medskip
  Ejemplo: \quad $\triop{$x$} = 2x + 3$
  \par\smallskip
  Luego: \quad $\triop{$2$} =$
  \par\smallskip
  \phantom{Luego:} \quad $\triop{$3$} =$

\item \rotulo{Mediante Tabla}
  En el conjunto $A = \{a, b, c, d\}$ podemos definir la siguiente tabla.

  \begin{cayley}{4}
    $\op$ & $a$ & $b$ & $c$ & $d$ \\ \cline{2-5}
    $a$   & $b$ & $c$ & $d$ & $a$ \\
    $b$   & $c$ & $d$ & $a$ & $b$ \\
    $c$   & $d$ & $a$ & $b$ & $c$ \\
    $d$   & $a$ & $b$ & $c$ & $d$ \\
  \end{cayley}

  Entonces: \quad $a \op b =$ \hfill $b \op c =$ \par\smallskip
  \phantom{Entonces:} \quad $a \op d =$ \hfill $c \op d =$

  \par\medskip
  Se puede usar cualquier símbolo para mi ``nueva operación matemática''.
  \par\smallskip
  Ejemplo: $\op$, \#, $\triangle$, $\theta$, $\bullet$, $\sim$, \dots, etc.

\end{apartados}

% ------------------------------ PROPIEDADES ------------------------------
\concepto{Propiedades de las Operaciones Matemáticas}

\begin{apartados}

\item \rotulo{Clausura o Cerradura}
  Si $a$ y $b$ pertenecen a un conjunto ``$C$'', por ejemplo, la operación
  definida también pertenece a dicho conjunto.

  \par\smallskip
  Ejemplo: En $\mathbb{N}$ la suma es cerrada:
  \[ 3 + 4 = 7 \]
  $3 \in \mathbb{N}$, $4 \in \mathbb{N}$ entonces $7 \in \mathbb{N}$

  \par\smallskip
  --- En $\mathbb{N}$ la multiplicación es cerrada:
  \[ 8 \times 5 = 40 \]
  $8 \in \mathbb{N}$, $5 \in \mathbb{N}$ entonces $40 \in \mathbb{N}$

  \par\smallskip
  \uline{Aplicación}:
  \par\smallskip
  --- En $\mathbb{N}$ se define: $a \op b = 3a + 4b$. ¿Es cerrada?

  \par\smallskip
  \uline{Solución}:
  \lineas{2}

  --- En $A = \{a, b, c\}$ se define la tabla:

  \begin{cayley}{3}
    $\op$ & $c$ & $b$ & $a$ \\ \cline{2-4}
    $a$   & $a$ & $b$ & $c$ \\
    $b$   & $b$ & $c$ & $a$ \\
    $c$   & $c$ & $a$ & $b$ \\
  \end{cayley}
  ¿Es cerrada?
  \lineas{2}

\item \rotulo{Propiedad Conmutativa}
  \formula{$\forall\, a, b \in C \qquad a \op b = b \op a$}

  Ejemplos: En $\mathbb{N}$ la suma es conmutativa.
  \[ 8 + 3 = 3 + 8 \qquad 2 + 7 = 7 + 2 \]

  --- En $\mathbb{Z}$ la multiplicación es conmutativa.
  \[ 8 \times 3 = 3 \times 8 \qquad 7 \times 2 = 2 \times 7 \]

  \uline{Aplicación}:
  \par\smallskip
  --- En $\mathbb{N}$ se define $a \,\theta\, b = a + b + 3$. ¿Será conmutativa?

  \par\smallskip
  \uline{Solución}:
  \lineas{2}

  --- En $C = \{m, n, p\}$ se define la tabla. ¿Es conmutativa?

  \begin{cayley}{3}
    $\theta$ & $m$ & $n$ & $p$ \\ \cline{2-4}
    $m$      & $n$ & $p$ & $m$ \\
    $n$      & $p$ & $m$ & $n$ \\
    $p$      & $m$ & $n$ & $p$ \\
  \end{cayley}

\item \rotulo{Elemento Neutro}
  Es aquel que, operando con cualquier número, se obtiene el mismo número.

  \par\smallskip
  Ejemplos:
  \par\smallskip
  --- El elemento neutro de la suma es el $0$.
  \[ 3 + 0 = 3 \,,\qquad 11 + 0 = 11 \]
  --- El elemento neutro de la multiplicación es el $1$.
  \[ 4 \times 1 = 4 \,,\qquad 19 \times 1 = 19 \ \text{etc.} \]

  \uline{Aplicación}:
  \par\smallskip
  --- En $\mathbb{N}$ se define: $a \,\triangle\, b = a - b + 2$.
  ¿Cuál será el elemento neutro?

  \par\smallskip
  \uline{Solución}:
  \lineas{2}

  --- En $B = \{x, y, z\}$ se define la tabla. ¿Cuál será el elemento neutro?

  \begin{cayley}{3}
    $\triangle$ & $x$ & $y$ & $z$ \\ \cline{2-4}
    $z$         & $x$ & $y$ & $z$ \\
    $x$         & $y$ & $z$ & $x$ \\
    $y$         & $z$ & $x$ & $y$ \\
  \end{cayley}

\item \rotulo{Elemento Inverso}
  Es aquel que, operando con un número, se obtiene el elemento neutro.
  El inverso de un número es único para ese número.

  \par\smallskip
  Ejemplo: En la suma el inverso de $4$ es $-4$, porque $4 + (-4) = 0$.

  \par\smallskip
  \uline{Aplicación}:
  \par\smallskip
  $*$ Del ejemplo anterior, para la operación $\triangle$ hallar el inverso
  de $3$ y el inverso de $5$.
  \par\smallskip
  Inverso de $3$ \;$(3^{-1}) =$ \par\smallskip
  Inverso de $5$ \;$(5^{-1}) =$

  \par\medskip
  $*$ Del ejemplo anterior de la tabla, hallar:
  \par\smallskip
  inverso de $x$ \;$(x^{-1}) =$ \par\smallskip
  inverso de $y$ \;$(y^{-1}) =$ \par\smallskip
  inverso de $z$ \;$(z^{-1}) =$

\end{apartados}

% ======================= EJERCICIOS DE APLICACIÓN =========================
\bloque{Ejercicios de Aplicación}

\begin{ejercicios}

\item La operación matemática es un proceso que consiste en la
  \rule{32mm}{0.4pt} de una o más \rule{22mm}{0.4pt} en otra cantidad
  llamada \rule{30mm}{0.4pt}.

\item La operación matemática es representada por un símbolo llamado
  \rule{34mm}{0.4pt}.

\item Si: $a \op b = 2a + b$. \; Hallar: $3 \op 4$
  \alts{9}{10}{11}{12}{13}

\item Se define en $A = \{a, b, c, d\}$ la siguiente tabla:
  \begin{cayley}{4}
    $\op$ & $a$ & $b$ & $c$ & $d$ \\ \cline{2-5}
    $a$   & $b$ & $c$ & $d$ & $a$ \\
    $b$   & $c$ & $d$ & $a$ & $b$ \\
    $c$   & $d$ & $a$ & $b$ & $c$ \\
    $d$   & $a$ & $b$ & $c$ & $d$ \\
  \end{cayley}
  Hallar: $(b \op d) \op (a \op c)$
  \alts{$a$}{$b$}{$c$}{$d$}{$b$ y $d$}

\item Se define: \; $\triop{$x$} = x^{2} + 3x$
  \par\smallskip
  Hallar: \; $\triop{$4$} + \triop{$5$}$
  \alts{66}{67}{68}{69}{70}

\item Si: $m \,\#\, n = 2m + 3n$. \; Hallar: $(2 \,\#\, 3) \,\#\, (4 \,\#\, 2)$
  \alts{76}{77}{78}{79}{80}

\item Se define:
  \[ a \,\triangle\, b = \begin{cases} 2a - b\,; & a > b \\ a + b\,; & a < b \end{cases} \]
  Calcular: $P = (2 \,\triangle\, 1) \,\triangle\, (1 \,\triangle\, 2)$
  \alts{5}{6}{7}{8}{9}

\item Si: $\triop{$\substack{b\\ a\ c}$} = a^{2} - bc$
  \par\smallskip
  Hallar: $\triop{$\substack{3\\ 4\ 2}$} - \triop{$\substack{2\\ 3\ 1}$}$
  \alts{1}{2}{3}{4}{5}

\item Se define en $A = \{2, 3, 4\}$:
  \begin{cayleyfull}{3}
    $\op$ & $2$ & $3$ & $4$ \\ \hline
    $2$   & $4$ & $3$ & $2$ \\ \hline
    $3$   & $2$ & $4$ & $3$ \\ \hline
    $4$   & $3$ & $2$ & $4$ \\ \hline
  \end{cayleyfull}
  Calcular: $S = \dfrac{(3 \op 4) \op (2 \op 4)}{(2 \op 3) \op (3 \op 4)}$
  \alts{1}{2}{0,5}{0,2}{3}

\item Dada la siguiente tabla:
  \begin{cayley}{4}
    $\op$ & $a$ & $b$ & $c$ & $d$ \\ \cline{2-5}
    $a$   & $c$ & $d$ & $a$ & $b$ \\
    $b$   & $d$ & $a$ & $b$ & $c$ \\
    $c$   & $a$ & $b$ & $c$ & $d$ \\
    $d$   & $b$ & $c$ & $d$ & $a$ \\
  \end{cayley}
  Calcular: $M = \dfrac{(a \op c) \op (b \op d)}{(a \op b) \op (c \op d)}$
  \alts{$b$}{$a$}{$a/b$}{1}{$d$}

\item Se tiene la siguiente tabla:
  \begin{cayley}{5}
    $\op$ & $m$ & $n$ & $p$ & $q$ & $r$ \\ \cline{2-6}
    $m$   & $p$ & $q$ & $m$ & $n$ & $r$ \\
    $n$   & $q$ & $p$ & $n$ & $r$ & $m$ \\
    $p$   & $m$ & $n$ & $p$ & $q$ & $r$ \\
    $q$   & $n$ & $r$ & $q$ & $p$ & $m$ \\
    $r$   & $r$ & $m$ & $r$ & $n$ & $p$ \\
  \end{cayley}
  Hallar el elemento neutro.
  \alts{$m$}{$n$}{$p$}{$q$}{$r$}

\item Del ejercicio anterior, hallar:
  $(n^{-1} \op p^{-1}) \op (q^{-1} \op r^{-1})$
  \alts{$m$}{$n$}{$p$}{$q$}{$r$}

\item Si: \; $\triop{$x$} = 2(x - 1)$ \quad y \quad $\triop{$\nabla x$} = 3(x-1)$
  \par\smallskip
  Hallar $x$ en: \; $\triop{$\nabla x$} = \triop{$2$}$
  \alts{4/7}{7/3}{13/7}{13/6}{13/3}

\item En el conjunto $A = \{0, 1, 2, 3, 4\}$:
  \begin{cayley}{4}
    $\triangle$ & $0$ & $1$ & $2$ & $3$ \\ \cline{2-5}
    $0$ & $2$ & $3$ & $0$ & $1$ \\
    $1$ & $2$ & $3$ & $0$ & $1$ \\
    $2$ & $0$ & $1$ & $1$ & $1$ \\
    $3$ & $3$ & $2$ & $1$ & $0$ \\
  \end{cayley}
  \begin{cayley}{4}
    $\bullet$ & $1$ & $2$ & $3$ & $4$ \\ \cline{2-5}
    $1$ & $1$ & $1$ & $1$ & $1$ \\
    $2$ & $2$ & $4$ & $1$ & $2$ \\
    $3$ & $1$ & $1$ & $4$ & $2$ \\
    $4$ & $1$ & $2$ & $2$ & $4$ \\
  \end{cayley}
  Hallar ``$x$'' en:
  $(x \,\triangle\, x) \bullet (3 \,\triangle\, 1) = (4 \,\triangle\, 3) \bullet (4 \,\triangle\, 1)$
  \alts{0}{1}{2}{3}{4}

\item Se define en $C = \{a, b, c, d, e\}$:
  \begin{cayley}{5}
    $\op$ & $a$ & $b$ & $c$ & $d$ & $e$ \\ \cline{2-6}
    $a$ & $c$ & $d$ & $e$ & $a$ & $b$ \\
    $b$ & $d$ & $e$ & $a$ & $b$ & $c$ \\
    $c$ & $e$ & $a$ & $b$ & $c$ & $d$ \\
    $d$ & $a$ & $b$ & $c$ & $d$ & $e$ \\
    $e$ & $b$ & $c$ & $d$ & $e$ & $a$ \\
  \end{cayley}
  ¿Cuál de las siguientes proposiciones es verdadera?
  \begin{enumerate}[label=\Roman*.,leftmargin=8mm,itemsep=0.2ex,topsep=0.4ex]
    \item $[a \op (x \op c)] \op d = c$, si $x = e$
    \item Se cumple la propiedad conmutativa
    \item Se cumple la propiedad de clausura
    \item El elemento neutro es ``$a$''
  \end{enumerate}
  \alts{I y III}{III y IV}{II y III}{I y IV}{Todas}

\end{ejercicios}

% ========================= TAREA DOMICILIARIA =============================
\tarea{5}

\begin{ejercicios}

\item Colocar verdadero (V) o falso (F) según corresponda:
  \begin{itemize}[label=--,leftmargin=6mm,itemsep=0.6ex,topsep=0.6ex]
    \item Toda operación matemática presenta una regla de definición. \hfill (\quad)
    \item El elemento neutro es aquél que operado con otro elemento se
          obtiene el elemento inverso. \hfill (\quad)
    \item La operación matemática es representada por el operador. \hfill (\quad)
    \item Toda operación matemática presenta elemento neutro. \hfill (\quad)
  \end{itemize}

\item El resultado de operar un elemento con su inverso es el
  \rule{34mm}{0.4pt}.

\item Si: $a \op b = a - 2b$. \; Hallar: $5 \op 2$
  \alts{1}{2}{3}{4}{5}

\item Se define en $A = \{a, b, c, d\}$ la siguiente tabla:
  \begin{cayley}{4}
    $\op$ & $a$ & $b$ & $c$ & $d$ \\ \cline{2-5}
    $a$   & $c$ & $d$ & $a$ & $b$ \\
    $b$   & $d$ & $a$ & $b$ & $c$ \\
    $c$   & $a$ & $b$ & $c$ & $d$ \\
    $d$   & $b$ & $c$ & $d$ & $a$ \\
  \end{cayley}
  Hallar: $(c \op a) \op (d \op b)$
  \alts{$a$}{$b$}{$c$}{$d$}{$a$ ó $c$}

\item Se define: \; $\triop{$x$} = 2x + 3$ \quad y \quad $\triop{$\nabla x$} = 4x - 5$
  \par\smallskip
  Hallar: \; $\triop{$5$} + \triop{$\nabla 3$}$
  \alts{19}{20}{21}{22}{23}

\item Si: $m \,\%\, n = 2m - n$ \; y \; $m \,\triangle\, n = n - 3m$
  \par\smallskip
  Hallar: $\dfrac{(4 \,\%\, 3)}{(20 \,\triangle\, 5)}$
  \alts{2}{1}{0}{-1}{-2}

\item Se define:
  \[ a \op b = \begin{cases} a - 2b\,; & a > b \\ a - b\,; & a < b \end{cases} \]
  Calcular: $M = (5 \op 2) \op (4 \op 3)$
  \alts{-2}{-1}{0}{1}{2}

\item Si: $\triop{$\substack{a\\ b\ c}$} = bc - a^{2}$
  \par\smallskip
  Hallar: $\triop{$\substack{2\\ 5\ 2}$} - \triop{$\substack{3\\ 2\ 7}$}$
  \alts{1}{2}{3}{4}{5}

\item Se define en $A = \{5, 6, 7\}$:
  \begin{cayleyfull}{3}
    $\#$ & $5$ & $6$ & $7$ \\ \hline
    $5$  & $7$ & $6$ & $5$ \\ \hline
    $6$  & $5$ & $7$ & $6$ \\ \hline
    $7$  & $6$ & $5$ & $7$ \\ \hline
  \end{cayleyfull}
  Calcular: $E = \dfrac{(6 \,\#\, 7) \,\#\, (5 \,\#\, 7)}{(5 \,\#\, 6) \,\#\, (6 \,\#\, 7)}$
  \alts{1}{2}{0,7}{0,2}{3}

\item Dada la siguiente tabla:
  \begin{cayley}{4}
    $\%$ & $d$ & $c$ & $d$ & $a$ \\ \cline{2-5}
    $a$  & $c$ & $d$ & $a$ & $b$ \\
    $b$  & $d$ & $a$ & $b$ & $c$ \\
    $c$  & $a$ & $b$ & $c$ & $d$ \\
    $d$  & $b$ & $c$ & $d$ & $a$ \\
  \end{cayley}
  Calcular: $N = \dfrac{(d \,\%\, b) \,\%\, (c \,\%\, a)}{(d \,\%\, c) \,\%\, (b \,\%\, a)}$
  \alts{$a$}{$b$}{$b/a$}{1}{$c$}

\item Se tiene la siguiente tabla:
  \begin{cayley}{5}
    $\op$ & $v$ & $w$ & $x$ & $y$ & $z$ \\ \cline{2-6}
    $z$ & $x$ & $y$ & $z$ & $v$ & $w$ \\
    $v$ & $y$ & $z$ & $v$ & $w$ & $x$ \\
    $w$ & $z$ & $v$ & $w$ & $x$ & $y$ \\
    $x$ & $v$ & $w$ & $x$ & $y$ & $z$ \\
    $y$ & $w$ & $x$ & $y$ & $z$ & $v$ \\
  \end{cayley}
  Hallar el elemento neutro.
  \alts{$v$}{$w$}{$x$}{$y$}{$z$}

\item Del ejercicio anterior, hallar:
  $[(w^{-1} \op z^{-1})] \op (y^{-1} \op v^{-1})] \op x^{-1}$

\item Si: \; $\triop{$x$} = x + 4$ \quad y \quad $\triop{$\nabla x$} = 5 - x$
  \par\smallskip
  Hallar ``$x$'' en: \; $\triop{$\nabla x$} = \triop{$3$}$
  \alts{9}{10}{11}{12}{13}

\item En el conjunto $B = \{0, 1, 2, 3, 4\}$:
  \begin{cayley}{4}
    $\triangle$ & $1$ & $2$ & $3$ & $4$ \\ \cline{2-5}
    $4$ & $2$ & $3$ & $4$ & $1$ \\
    $3$ & $3$ & $4$ & $1$ & $2$ \\
    $2$ & $4$ & $1$ & $2$ & $3$ \\
    $1$ & $1$ & $2$ & $3$ & $4$ \\
  \end{cayley}
  \begin{cayley}{4}
    $\bullet$ & $3$ & $2$ & $1$ & $0$ \\ \cline{2-5}
    $0$ & $3$ & $2$ & $1$ & $0$ \\
    $1$ & $2$ & $1$ & $0$ & $3$ \\
    $2$ & $1$ & $0$ & $3$ & $2$ \\
    $3$ & $0$ & $3$ & $2$ & $1$ \\
  \end{cayley}
  Hallar ``$x$'' en:
  $(3 \,\triangle\, x) \bullet (4 \,\triangle\, 1) = (3 \bullet 2) \,\triangle\, (1 \bullet 0)$
  \alts{0}{1}{2}{3}{4}

\item Se define en $C = \{m, n, p, q, r\}$:
  \begin{cayley}{5}
    $\op$ & $m$ & $n$ & $p$ & $q$ & $r$ \\ \cline{2-6}
    $r$ & $p$ & $q$ & $r$ & $m$ & $n$ \\
    $q$ & $q$ & $r$ & $m$ & $n$ & $p$ \\
    $p$ & $r$ & $m$ & $n$ & $p$ & $q$ \\
    $n$ & $m$ & $n$ & $p$ & $q$ & $r$ \\
    $m$ & $r$ & $p$ & $q$ & $r$ & $m$ \\
  \end{cayley}
  ¿Cuál de las siguientes proposiciones es verdadera?
  \begin{enumerate}[label=\Roman*.,leftmargin=8mm,itemsep=0.2ex,topsep=0.4ex]
    \item Es conmutativa
    \item Es cerrada
    \item No tiene elemento neutro
  \end{enumerate}
  \alts{I}{II y III}{III}{I y II}{Todas}

\end{ejercicios}

\end{ficha}
\end{document}
